% ===========================================================================
% Section 2. Conceptual framework (input by main.tex; no preamble).
% Framing: NEIS working paper on the Allied Industrial Base (neis2026allied): definition, the
% Severity -> Coordination -> Specialization logic and the five analytical layers (its Sections
% 2, 7 and 8), paraphrased and cited.
% Specification (code is the specification): config/weights.yml, config/rankability.yml,
%   R/utils/scurve.R, R/indices/index_builder_core.R, R/indices/build_policy_index.R,
%   R/indices/build_indices.R (Strategic Index, endowment screen, allied pick),
%   R/geopolitical_distance/*.R and scripts/12 (alignment), R/indices/coalition.R and scripts/13
%   (tiers, EU node), R/themes/partnership_strength/*.R (FS, PO, DP, PSI),
%   R/indices/allied_network_design.R and scripts/30 (network; lambda = 0.78, node policy term off).
% Figure fig:framework_architecture is drawn inline (TikZ) so that it matches the final method.
% ===========================================================================
\section{Conceptual framework}\label{sec:framework}

This section places the paper within the analytical architecture proposed for an Allied
Industrial Base \citep{neis2026allied}, states which parts of that architecture the paper
measures and which it leaves aside, and then defines the objects used in the rest of the paper.
Figure~\ref{fig:framework_architecture} summarises the construction and Table~\ref{tab:notation}
collects the notation.

% ---------------------------------------------------------------------------
\subsection{The allied industrial base as a design problem}\label{subsec:framework_aib}

The working paper defines an Allied Industrial Base by its purpose: trusted economies coordinate
strategically important capabilities so that, together, they reach a resilience and scale that
none could reach efficiently alone \citep{neis2026allied}. The object is
the network: a capability should exist at sufficient scale somewhere among trusted partners, not
in every member. The unit of policy is the \emph{node}, a step such as mineral refining or
component manufacture, because one supply chain can warrant different treatment at different
steps. And the problem is one of allocation: which capabilities should be domestic, duplicated,
specialised among partners or left to global markets.

The working paper orders the decision in three steps. \emph{Severity} asks how damaging a
disruption would be and how hard substitution or recovery would be; \emph{coordination}, whether
the problem needs international action rather than private diversification or national policy;
and \emph{specialization}, where the capability should reside and what should be duplicated. The
evidence is organised in five analytical layers: node criticality and severity, country
capability and opportunity, partner complementarity and reciprocity, systemic bankability, and
network robustness under stress. The object is a combination of capabilities, not a ranking: a
partner that ranks low in isolation can be indispensable if it supplies what the network lacks.

Table~\ref{tab:framework_layers} maps the layers onto this paper. We measure each economy's
relative exposure within a node by the need term $1-\ES$ (rank-based, so it does not compare
severity across nodes) and, for extraction, apply the endowment screen; capability by
$\EO$, with the Policy index recording where governments already concentrate their effort;
complementarity, alignment and part of reciprocity by the coalition and the bilateral indices;
and robustness only in part, by the diversification penalty and multi-sourcing constraints of the
network design. Systemic bankability \citep{spokas2025systemic}, common-mode risk (beyond a first
measure of the Chinese ownership of suppliers' foreign-invested plants), dynamic option value,
reciprocal obligations, resilience floors and stress tests are not modelled.

\begin{table}[tbp]
\centering
\begin{threeparttable}
\caption{The analytical layers of an Allied Industrial Base and their counterparts in this paper}
\label{tab:framework_layers}
\footnotesize
\begin{tabularx}{\textwidth}{@{}>{\raggedright\arraybackslash}p{2.7cm}
  >{\raggedright\arraybackslash}X >{\raggedright\arraybackslash}X@{}}
\toprule
Layer & Measured here (section) & Not measured\\
\midrule
1. Node criticality and severity &
  Country-relative exposure within each technology--stage pair, $1-\ES_{its}$ (import
  dependence, supplier concentration, reserves; \ref{sec:es}); endowment screen on extraction
  stages (\ref{sec:strategic}) &
  Severity compared across nodes; disruption horizons, inventories, qualification and build times, end-use criticality\\[2pt]
2. Country capability and opportunity &
  $\EO_{its}$ by cell (\ref{sec:eo}); Policy index and Strategic Index (\ref{sec:strategic}) &
  Plant-level cost curves, pipeline capacity, workforce, tacit know-how\\[2pt]
3. Complementarity and reciprocity &
  Coalition rule from UN-vote alignment, agreements and supply shares (\ref{sec:coalition});
  Friendshore, Partnership Opportunity and Development Potential (\ref{sec:bilateral}); network
  edges (\ref{sec:network}) &
  Reciprocal obligations, export-restriction risk, market-access commitments\\[2pt]
4. Systemic bankability & --- &
  Offtake, price risk, permitting, enabling infrastructure, insurance, delivery capacity\\[2pt]
5. Network robustness &
  Concave penalty on the Herfindahl index of designed shares, share bounds and multi-sourcing
  constraints, calibrated on observed concentration (\ref{sec:network}); Chinese-parent exposure in
  four Midstream cells (\ref{sec:bilateral}) &
  Stress tests, common-mode failure through shared inputs, logistics or equipment, stockpiles,
  recovery time, dynamic option value\\
\bottomrule
\end{tabularx}
\begin{tablenotes}[flushleft]\footnotesize
\item \textit{Notes:} Layers and their content follow \citet{neis2026allied}; numbers in
parentheses refer to sections of this paper.
\end{tablenotes}
\end{threeparttable}
\end{table}

Two simplifications matter. First, our cell, a technology--stage pair, is coarser than the
working paper's node: Batteries Midstream, for example, groups mineral refining, electrode
materials and cell manufacture. It is the finest level at which both pillars can be scored for
most economies from public data. Second, the coordination step is assumed rather than tested: the
22 traded Upstream and Midstream stages are all treated as candidates, while Downstream
deployment, largely non-traded, enters the Strategic Index but not the network design.

% ---------------------------------------------------------------------------
\subsection{Cells, pillars and rankability}\label{subsec:framework_unit}

The unit of analysis is the cell $\cell$: an economy $i\in\Countries$, a technology
$t\in\Techs$ and a stage $s\in\Stages$, where Upstream is the extraction of fuels and minerals,
Midstream their processing and the manufacture of components and equipment, and Downstream
deployment and end use. $\Techs$ holds the eleven headline technologies, fossil fuels included,
so an economy has at most 33 cells, and $\Countries$ holds 219 economies after harmonisation,
195 sovereign states and 24 territories (Section~\ref{sec:data}).

Energy Security $\ES_{its}$ measures how secure $i$'s supply of technology $t$ at stage $s$ is,
higher values being more secure. It combines exposure (import dependence, the concentration of
supplying markets) with resilience (reserves, production, consumption, access and investment).
Economic Opportunity $\EO_{its}$ measures how well placed $i$ is to capture value as a producer,
from comparative advantage, export feasibility, demand, cost, technological readiness and
investment. Security describes an economy as a user and opportunity as a producer, so the two can
diverge: a fuel importer with a strong manufacturing base can be insecure Upstream and well
placed Midstream. Shared components make part of their association mechanical
(Section~\ref{sec:joint}).

Both pillars are built in the same way (Section~\ref{sec:normalization}). Each raw variable
$x_{itsv}$, oriented so that higher is better, is replaced by the S-curve of its percent rank
within a comparison set $g$, usually the economies observed in the pair $(t,s)$,
\begin{equation}\label{eq:fw_scurve}
  \tilde x_{itsv}=\Scurve\bigl(r_{itsv}\bigr),\qquad
  \Scurve(r)=\frac{r^{\gamma}}{r^{\gamma}+(1-r)^{\gamma}},\qquad
  r_{itsv}=\frac{\rank_g\bigl(x_{itsv}\bigr)-1}{n_g-1},\qquad \gamma=0.5 ,
\end{equation}
the probability-weighting function of \citet{karmarkar1978subjectively}, which stretches the tails
of a distribution relative to its middle. A category averages the components observed, or filled
by an explicit missing-data rule, for the cell,
\begin{equation}\label{eq:fw_category}
  c_{itsk}=\frac{1}{|\mathcal{V}^{*}_{itsk}|}\sum_{v\in\mathcal{V}^{*}_{itsk}}\hat x_{itsv},
\end{equation}
and the pillar is the weighted mean over the categories $\Cats^{P}_{its}$ observed for the cell,
\begin{equation}\label{eq:fw_pillar}
  P_{its}=\frac{\sum_{k\in\Cats^{P}_{its}}w^{P}_{k}\,c_{itsk}}{\sum_{k\in\Cats^{P}_{its}}w^{P}_{k}},
  \qquad P\in\{\ES,\EO\},
\end{equation}
with ten categories and fixed weights per pillar (2 to 6 for $\ES$, 2 to 4 for $\EO$). An
unobserved category drops out with its weight rather than counting as zero, and each pillar
decomposes exactly into category and variable contributions. A cell enters a ranking, a top-$N$
list or the Strategic Index only if it is \emph{rankable}: the economy is sovereign and at most
half of the cell's index weight is imputed (equation~\eqref{eq:rankable}). This holds for 4,163
of 6,232 scored $\ES$ cells and 4,301 of 7,131 $\EO$ cells.

% ---------------------------------------------------------------------------
\subsection{Policy and the Strategic Index}\label{subsec:framework_si}

The Policy index $\POL_{its}$ records where a government directs its policy. It is the weighted
mean of an index of the measures in the IEA Policies and Measures database tagged to the cell
(weight 3, required) and a component that interacts the country's Climate Action Tracker rating
with the technology's emission intensity (weight 2, where the country is rated). It covers 215
economies, 65 on both components and 150 on the IEA component alone, and the basis matters
(Section~\ref{sec:strategic}). Measures from the New Industrial Policy Observatory enter only as a
robustness variant (Section~\ref{sec:robustness}). The Strategic Index ranks the sectors of each
country by the case for building there,
\begin{equation}\label{eq:fw_si}
  \SI_{its}=0.2\,\POL^{*}_{its}+0.3\,\EO_{its}+0.3\,\bigl(1-\ES_{its}\bigr)
  +0.1\,\lambda_s+0.1\,g_t ,
\end{equation}
where $\POL^{*}_{its}$ is the S-curve of $\POL$ across $i$'s own scored sectors, the pillars enter
at their cross-country levels, $\lambda_s$ is a stage weight (Midstream 0.75, Upstream 0.5,
Downstream 0.25) and $g_t\in[0,1]$ rises as lifecycle greenhouse-gas intensity falls. A cell is
scored only if the country has a Policy index and both pillar cells are rankable; nothing is
imputed.

Security enters as \emph{need}, for a reason of marginal value. Let $V_i$ be the value to $i$ of
supply security in a sector, increasing and concave, and $\rho_i$ the probability that an effort
to build capacity there succeeds, increasing in opportunity and in existing policy commitment. An
effort that raises security by $\Delta$ is worth about
\begin{equation}\label{eq:fw_marginal}
  \rho_i\bigl(\EO_{its},\POL_{its}\bigr)\,V_i'\bigl(\ES_{its}\bigr)\,\Delta ,
\end{equation}
which rises with opportunity and policy and falls with the security already attained. The index
linearises this ranking, so sectors in which a country is exposed but capable rank first; reading
security as strength would rank existing strengths first, which $\EO$ already describes.
Equation~\eqref{eq:fw_marginal} motivates the signs, not the additive form or the weights, which
are normative \citep{decancq2013weights}; Section~\ref{sec:robustness} reports the strength
reading.

Each coalition member's \emph{allied pick} is its highest-scoring Upstream or Midstream sector
outside geothermal and green hydrogen. Read as need, $\ES$ gives an economy without a resource
its largest need at the extraction stage, so an \emph{endowment screen} admits an Upstream stage
to the pick only if the member has recorded reserves or production of the stage's resources. The
screen is a selection rule, not part of the index: values are unchanged, and 669 of the 998 scored
Upstream cells lack a resource base. The network design uses neither the index nor the picks.

% ---------------------------------------------------------------------------
\subsection{Coalition, bilateral indices and network design}\label{subsec:framework_layers}

\paragraph{The coalition.} The coalition $\Coalition$ is derived by rule
(Section~\ref{sec:coalition}). Each voting economy is placed on a U.S.--China axis
$s_i\in[0,1]$ from its UN General Assembly ideal point (equation~\eqref{eq:coalition_axis}), and
quartiles of $s_i$ define U.S.-leaning, nonaligned and China-leaning blocs. The 31 \emph{core}
members are the U.S.-leaning economies with GDP of at least USD 100 billion. The 26
\emph{partners} are members of a U.S. free-trade agreement, the Indo-Pacific Economic Framework,
the Minerals Security Partnership, the G7 or NATO above the same GDP floor, and major suppliers
outside the China-leaning bloc without a GDP floor: at least 5 percent of world production,
refining, manufacturing capacity or exports in an Upstream or Midstream stage, or at least 10
percent of world reserves of a headline mineral, fuel or uranium. With the six EU members below
the GDP floor the coalition spans 63 economies; the network treats the EU-27 as one node with
stage-weighted member inputs, which gives 37 nodes. The dyadic alignment $A_{ij}$
(equation~\eqref{eq:coalition_dyadic}), the probability that two economies fall on the same side
of the divide, links the coalition to the bilateral layer.

\paragraph{Bilateral indices.} For a reporter $i$, a partner $j\neq i$ and a pair $(t,s)$,
Friendshore $\FS_{ijts}$ rates $j$ as a supplier to $i$ from $i$'s imports from $j$, $i$'s need,
$j$'s opportunity, $i$'s outward investment in $j$, the alignment $A_{ij}$ and, in four
manufacturing stages, the share of $j$'s foreign-invested plants without Chinese parents, less a
climate-policy penalty (equation~\eqref{eq:fs_index}). Partnership Opportunity $\PO_{ijts}$ rates
$j$ as a market for $i$'s exports, and Development Potential $\DP_{jts}$ scores partners that are
not high-income economies. The Partnership Strength Index combines best-five summaries of the
three with weights 0.4, 0.4 and 0.2, those of the partnership score of \citet{carey2025secure};
the headline averages the views of the 31 core members. Table~\ref{tab:roles} shows how security
and opportunity enter each index.

\paragraph{Network design.} For each of the 22 Upstream and Midstream stages, a mixed-integer
program over the 37 nodes chooses producers ($y_p\in\{0,1\}$), supplier links
($\ell_{pc}\in\{0,1\}$) and flows $f_{pc}\ge0$, the share of the network's unit demand that $p$
supplies to $c$, with production shares $x_p=\sum_c f_{pc}$:
\begin{equation}\label{eq:fw_network}
  \max_{y,\ell,f}\ \sum_{p}\sum_{c}f_{pc}\bigl(a_p+0.5\,W_{pc}\bigr)-\lambda\sum_{p}x_p^{2}
  \quad\text{s.t.}\quad \sum_{p}f_{pc}=d_c,\quad 0.05\,y_p\le x_p\le0.40\,y_p .
\end{equation}
The producer score $a_p$ weights the S-curves of $\EO$ and $\ES$ across the nodes by $5/7$ and
$2/7$, without a policy term; the edge weight $W_{pc}=0.5\,\FS_{cpts}+0.5\,\PO_{pcts}$ reads each
directed index from the end of the relationship it describes; and the demand weight $d_c$ mixes
the consumer's need and economic size. At least three nodes produce, and every consumer draws on
at least two suppliers, each providing at least 10 percent of its demand. The penalty on the
Herfindahl index of production shares makes the program spread production among nodes rather than
rank them, a first step toward optimising combinations. Because the penalty treats each node as an
independent source, it does not register that producers sharing a bloc, a region or a parent company
diversify less; Section~\ref{sec:discussion} discusses a correlated-risk version. Its weight is calibrated:
$\lambda=0.78$ is the smallest value at which the median designed Herfindahl index across stages
does not exceed the median observed export concentration of the nodes, a knife-edge choice that
Section~\ref{sec:network} discloses. An outer loop caps, by the logarithm of GDP, the number of
stages in which a node is among the five largest producers.

\begin{figure}[tbp]
  \centering
  \resizebox{\textwidth}{!}{%
  \begin{tikzpicture}[
    font=\scriptsize,
    >={Stealth[length=1.8mm]},
    box/.style={draw=#1, fill=#1!10, thick, rounded corners=2pt, align=center, inner sep=3pt,
                anchor=north},
    flow/.style={->, draw=osiink!75, line width=0.6pt},
    bus/.style={draw=osiink!75, line width=0.9pt},
    lab/.style={font=\scriptsize, text=osiink, fill=white, inner sep=1.2pt, align=center}
  ]
  \node[box=osigray, text width=15.4cm] (src) at (8, 0) {%
    {\footnotesize\bfseries Sources} by economy, technology and stage: IEA, Energy Institute,
    EIA, USGS, OECD-NEA Red Book, UN Comtrade (HS6 baskets, all-reporter dyads), Atlas of
    Economic Complexity, BloombergNEF, World Bank, IMF, CM-FDI;\\ IEA Policies and Measures,
    Climate Action Tracker; UN General Assembly ideal points; trade and security agreements};
  \node[box=osiaqua, text width=3.5cm, minimum height=2.5cm] (pol)
    at ($(src.south)+(-6,-0.5)$) {%
    {\footnotesize\bfseries Policy}\\ $\POL_{its}$\\[1pt] IEA PAMS (3, required)\\
    $+$ CAT (2) where rated;\\ 215 economies};
  \node[box=osiblue, text width=3.5cm, minimum height=2.5cm] (es)
    at ($(src.south)+(-2,-0.5)$) {%
    {\footnotesize\bfseries Energy Security}\\ $\ES_{its}$\\[1pt] 10 categories;\\
    higher $=$ more secure;\\ rankable if $\iota_{its}\le0.5$};
  \node[box=osiorange, text width=3.5cm, minimum height=2.5cm] (eo)
    at ($(src.south)+(2,-0.5)$) {%
    {\footnotesize\bfseries Economic Opportunity}\\ $\EO_{its}$\\[1pt] 10 categories;\\
    rankable if $\iota_{its}\le0.5$};
  \node[box=osiviolet, text width=3.5cm, minimum height=2.5cm] (coal)
    at ($(src.south)+(6,-0.5)$) {%
    {\footnotesize\bfseries Coalition}\\ $\Coalition$\\[1pt] UN-vote alignment,\\
    agreements, supply shares;\\ 31 core $+$ 26 partners;\\ EU-27 as one: 37 nodes};
  \coordinate (busy) at ($(src.south)+(0,-3.35)$);
  \draw[bus] ($(busy)+(-7.7,0)$) -- ($(busy)+(7.7,0)$);
  \foreach \n in {pol, es, eo, coal} {\draw[bus] (\n.south) -- (\n.south |- busy);}
  \node[box=osiviolet, text width=6.5cm, minimum height=2.2cm] (si)
    at ($(busy)+(-4.2,-0.6)$) {%
    {\footnotesize\bfseries Strategic Index}\\[1pt]
    $\SI_{its}=0.2\,\POL^{*}_{its}+0.3\,\EO_{its}+0.3\,(1-\ES_{its})+0.1\,\lambda_s+0.1\,g_t$\\
    rankable cells only; endowment screen on Upstream picks\\
    $\Rightarrow$ each member's allied pick (reported)};
  \node[box=osiyellow, text width=6.5cm, minimum height=2.2cm] (bil)
    at ($(busy)+(4.2,-0.6)$) {%
    {\footnotesize\bfseries Bilateral indices}\\[1pt]
    $\FS_{ijts}$: $j$ as supplier to $i$ (with $A_{ij}$ and Chinese-parent exposure)\\
    $\PO_{ijts}$: $j$ as market for $i$;\quad $\DP_{jts}$: development partner\\
    $\Rightarrow$ $\PSI_j=0.4\,\FS+0.4\,\PO+0.2\,\DP$, core view (reported)};
  \draw[flow] (si.north |- busy) -- node[lab, pos=0.45] {$\POL$, $\ES$, $\EO$} (si.north);
  \draw[flow] (bil.north |- busy) -- node[lab, pos=0.45] {$\ES$, $\EO$, $A_{ij}$} (bil.north);
  \node[box=osigray, text width=15.4cm] (net) at ($(busy)+(0,-4.3)$) {%
    {\footnotesize\bfseries Allied network design}: each of 22 Upstream and Midstream stages over
    37 nodes,\quad $\max\ \sum_{p,c}f_{pc}\,(a_p+0.5\,W_{pc})-\lambda\sum_p x_p^2$\\[1pt]
    $a_p$ from $\EO$ and $\ES$ (no policy term);\quad $W_{pc}=0.5\,\FS_{cp}+0.5\,\PO_{pc}$;\quad
    $d_c$ from need and size;\quad shares in $\{0\}\cup[0.05,0.40]$;\quad $\ge2$ suppliers per
    consumer;\\ $\lambda=0.78$ calibrated on the nodes' observed export concentration;\quad
    GDP-log portfolio caps};
  \draw[flow] (busy) -- node[lab, pos=0.5] {$\ES$, $\EO$,\\ nodes} (net.north);
  \draw[flow] (bil.south) -- node[lab, pos=0.45] {$\FS$, $\PO$} (bil.south |- net.north);
  \end{tikzpicture}%
  }
  \caption{Architecture of the analysis. Sources feed four cell- or economy-level constructs:
  the Policy index, the two pillars, and the coalition, derived from alignment in the UN General
  Assembly, agreements and supply shares. The Strategic Index ranks each country's sectors and
  yields the allied picks; the bilateral indices rate partners as suppliers and markets and are
  summarised in the Partnership Strength Index. The network design allocates production among
  the coalition's nodes. In the terms of \citet{neis2026allied}, the pillars, the Policy index and
  the Strategic Index address relative exposure and capability, the coalition and the bilateral indices
  complementarity, and the network design specialisation and, in part, robustness. $\iota_{its}$
  is the imputed share of a cell's index weight; the Strategic Index and the Partnership Strength
  Index are not inputs to the network design.}
  \label{fig:framework_architecture}
\end{figure}

\begin{table}[tbp]
\centering
\begin{threeparttable}
\caption{How security and opportunity enter each index}
\label{tab:roles}
\footnotesize
\begin{tabularx}{\textwidth}{@{}>{\raggedright\arraybackslash}p{2.7cm}
  >{\raggedright\arraybackslash}p{2.4cm} >{\raggedright\arraybackslash}X
  >{\raggedright\arraybackslash}X@{}}
\toprule
Index & Whose position & Energy security enters as & Economic opportunity enters as\\
\midrule
Strategic Index $\SI_{its}$ & country $i$, deciding where to build &
  need of $i$: $1-\ES_{its}$ (level), weight 0.3 &
  capability of $i$: $\EO_{its}$ (level), weight 0.3\\[2pt]
Friendshore $\FS_{ijts}$ & importer $i$, supplier $j$ &
  need of the importer, ranked across $i$'s own cells, weight 2 of 11 &
  capability of the supplier, ranked across economies in $(t,s)$, weight 2 of 11\\[2pt]
Partnership Opportunity $\PO_{ijts}$ & exporter $i$, market $j$ &
  need of the market, ranked across markets, weight 1 of 7 &
  capability of the exporter, ranked across economies, weight 3 of 7\\[2pt]
Development Potential $\DP_{jts}$ & partner $j$, not high income &
  need of the partner, weight 1 of 10 & opportunity of the partner, weight 1.5 of 10\\[2pt]
Network producer score $a_p$ & node $p$ as producer &
  security of the producer, S-curve across nodes, weight 2/7 &
  capability of the producer, S-curve across nodes, weight 5/7\\[2pt]
Network demand $d_c$ & node $c$ as consumer &
  need of the consumer, weight 0.5 (with size, 0.5) & ---\\
\bottomrule
\end{tabularx}
\begin{tablenotes}[flushleft]\footnotesize
\item \textit{Notes:} Weights are nominal, before renormalisation over the components observed.
Security enters as need wherever it describes the party that would gain from new supply, and as
security only in the producer score of the network design, where a node's own security in a
stage proxies its reliability as a supplier. Sections~\ref{sec:strategic}, \ref{sec:bilateral}
and~\ref{sec:network} give the full definitions.
\end{tablenotes}
\end{threeparttable}
\end{table}

\begin{table}[tbp]
\centering
\begin{threeparttable}
\caption{Notation}
\label{tab:notation}
\footnotesize
\begin{tabularx}{\textwidth}{@{}>{\raggedright\arraybackslash}p{2.9cm}
  >{\raggedright\arraybackslash}X >{\raggedright\arraybackslash}p{2.7cm}@{}}
\toprule
Symbol & Meaning & Defined in\\
\midrule
\multicolumn{3}{@{}l}{\textit{Sets and indices}}\\
$\Countries$ & the 219 economies after harmonisation & Sec.~\ref{sec:data}\\
$\Techs$, $\Stages$ & eleven headline technologies; \{Upstream, Midstream, Downstream\} &
  Sec.~\ref{subsec:framework_unit}\\
$\Coalition$ & the 57 coalition members (31 core, 26 partners); 37 network nodes &
  Sec.~\ref{sec:coalition}\\
$i,j$; $t,s$; $k,v$ & economy (reporter), partner; technology, stage; category, variable & \\
$p,c$ & network node as producer; as consumer & Sec.~\ref{sec:network}\\
$\cell$, $\dyad$ & cell; directed dyad with reporter $i$ and partner $j\neq i$ &
  Secs.~\ref{subsec:framework_unit}, \ref{sec:bilateral}\\[2pt]
\multicolumn{3}{@{}l}{\textit{Normalisation and pillars}}\\
$x_{itsv}$, $r$ & raw variable; percent rank within comparison set $g$ & Eq.~\eqref{eq:fw_scurve}\\
$\Scurve$, $\gamma$ & S-curve; curvature, $\gamma=0.5$; ${\Scurve}[\cdot]_g$ ranks within $g$ &
  Eq.~\eqref{eq:fw_scurve}\\
$\tilde x_{itsv}$, $\hat x_{itsv}$ & theme index; after the missing-data rule &
  Eqs.~\eqref{eq:fw_scurve}, \eqref{eq:fw_category}\\
$c_{itsk}$, $w^{P}_k$ & category score; weight of category $k$ in pillar $P$ &
  Eqs.~\eqref{eq:fw_category}, \eqref{eq:fw_pillar}\\
$\ES_{its}$, $\EO_{its}$ & Energy Security and Economic Opportunity pillars, in $[0,1]$ &
  Eq.~\eqref{eq:fw_pillar}\\
$\iota_{its}$ & imputed share of the cell's index weight; rankable if $\iota_{its}\le0.5$ &
  Eq.~\eqref{eq:rankable}\\[2pt]
\multicolumn{3}{@{}l}{\textit{Policy and Strategic Index}}\\
$\POL_{its}$, $\POL^{*}_{its}$ & Policy index; its S-curve across $i$'s scored sectors &
  Eqs.~\eqref{eq:pol_composite}, \eqref{eq:si_pol}\\
$\lambda_s$, $g_t$ & stage weight; technology emissions weight (1 = cleanest) &
  Eqs.~\eqref{eq:fw_si}, \eqref{eq:ghg_index}\\
$\SI_{its}$ & Strategic Index & Eq.~\eqref{eq:fw_si}\\[2pt]
\multicolumn{3}{@{}l}{\textit{Alignment and bilateral indices}}\\
$s_i$, $A_{ij}$ & position on the U.S.--China axis; dyadic alignment &
  Eqs.~\eqref{eq:coalition_axis}, \eqref{eq:coalition_dyadic}\\
$\FS_{ijts}$, $\PO_{ijts}$, $\DP_{jts}$ & Friendshore; Partnership Opportunity; Development
  Potential & Eqs.~\eqref{eq:fs_index}, \eqref{eq:po_index}, \eqref{eq:dp_index}\\
$\PSI_j$ & Partnership Strength Index & Eq.~\eqref{eq:psi}\\[2pt]
\multicolumn{3}{@{}l}{\textit{Network design}}\\
$a_p$, $W_{pc}$, $d_c$ & producer score; edge weight; demand weight ($\sum_c d_c=1$) &
  Eq.~\eqref{eq:fw_network}\\
$y_p$, $\ell_{pc}$, $f_{pc}$, $x_p$ & producer selection; supplier link; flow share; production
  share & Eq.~\eqref{eq:fw_network}\\
$\lambda$, $\HHI$ & diversification weight; Herfindahl--Hirschman index & Eq.~\eqref{eq:fw_network}\\
\bottomrule
\end{tabularx}
\begin{tablenotes}[flushleft]\footnotesize
\item \textit{Notes:} The stage weight $\lambda_s$ carries a subscript and the diversification
weight $\lambda$ none; the axis position $s_i$ carries an economy subscript and the stage $s$
none. The production share $x_p$ and the policy components $x^{v}_{its}$ of
Section~\ref{sec:strategic} are distinct from the raw variable $x_{itsv}$. Later sections
introduce further symbols where they are needed.
\end{tablenotes}
\end{threeparttable}
\end{table}
